% ═════════════════════════════════════════════════════════════════════════════
\section{Monotonic Abstraction}
\label{sec:monotonic}
% ═════════════════════════════════════════════════════════════════════════════

Regular model checking is a uniform framework in which a wide class of
techniques can be explained. In practice, however, it often needs to be
complemented by other techniques. A limiting factor is the computation of the
transitive closure of a transducer, which is not guaranteed to terminate and
may rely on expensive automata-theoretic constructions. One way to increase
efficiency is to over-approximate the behavior of the system. We derive an
\emph{abstract} system whose behavior includes that of the original,
\emph{concrete}, system, and that is easier to analyze. The price is a loss of
precision: the abstract system may contain \emph{false positives}, that is,
violations of the safety property that do not exist in the concrete system.
The method is nevertheless sound: if the abstract system satisfies a safety
property, then so does the concrete system.

In this section, we present an over-approximation, called \emph{monotonic
abstraction}~\cite{ADHR07,ABDR08}, that produces a well-structured transition
system, so
that backward reachability analysis can be applied
to the abstract system. We consider the parameterized systems of
Section~\ref{sec:prelim}, in which the transitions of a process may be
constrained by universal or existential global conditions. Global conditions are important in the description of
many parameterized systems, such as mutual exclusion protocols, but they are
difficult to analyze: parameterized systems with universal conditions can
simulate counter machines, so all non-trivial verification problems for them
are undecidable. Such systems can be described with transducers as in
Section~\ref{sec:rmc}, but here we use over-approximation instead.

\subsection{Example}
\label{sec:mutex}

We introduce the method through a protocol that implements mutual exclusion
among an arbitrary number of processes. Each process, depicted in
Figure~\ref{fig:mutex}, has four local states: idle~($\Idle$),
requesting~($\Req$), waiting~($\Wait$), and critical~($\Crit$). Initially, all
processes are idle. When a process becomes interested in the critical section,
it declares its interest by moving to the requesting state. This is described
by the universal transition~$t_1$, which is allowed only if all other
processes, whether in the left or the right context, are idle or requesting;
this is indicated by the index $LR$ of the quantifier. From the requesting
state, a process may move to the waiting state by the local transition~$t_2$,
which does not depend on the other processes. Any number of processes may
move from idle to requesting. However, once the first process has moved to the
waiting state, it closes the door on the processes that are still idle: they
cannot leave the idle state until no process is waiting or critical.

Among the processes that have declared their interest, the leftmost process
has the highest priority. This is described by the universal
transition~$t_4$, by which a process may move from the waiting state to the
critical section only if all processes in its left context are idle; the
index~$L$ stands for left. If there is a process in its left context that is
requesting, waiting, or critical, the process loops in the waiting state by the
existential transition~$t_3$. When the process leaves the critical section, it
returns to the requesting state by the local transition~$t_5$. From the
requesting state, it either tries to reach the critical section again, or
becomes idle by the local transition~$t_6$.

\begin{figure}[t]
\centering
\begin{tikzpicture}[->,>={Stealth[round]},shorten >=1pt,semithick,
    every state/.style={minimum size=7mm,inner sep=1pt}]
  \node[state,initial,initial text={},initial where=above] (i) at (0,0) {$\Idle$};
  \node[state] (r) at (0,-1.8) {$\Req$};
  \node[state] (w) at (-1.6,-3.6) {$\Wait$};
  \node[state] (c) at (1.6,-3.6) {$\Crit$};
  \path (i) edge[bend right=30] node[left,font=\small] {$t_1\colon \forall_{LR}\{\Idle,\Req\}$} (r)
        (r) edge[bend right=30] node[right,font=\small] {$t_6$} (i)
        (r) edge node[left,font=\small] {$t_2$} (w)
        (c) edge node[right,font=\small] {$t_5$} (r)
        (w) edge[loop left] node[left,font=\small] {$t_3\colon \exists_{L}\{\Req,\Wait,\Crit\}$} (w)
        (w) edge node[below,font=\small] {$t_4\colon \forall_{L}\{\Idle\}$} (c);
\end{tikzpicture}
\caption{One process in a mutual exclusion protocol with a linear topology.
The labels give the global conditions of the transitions $t_1$, $t_3$, and
$t_4$; the other transitions are local.}
\label{fig:mutex}
\end{figure}

\subsection{Model}
\lukas{introduce somewhere the indexed quantifiers used in the figure}
In the notation of Section~\ref{sec:prelim}, the protocol is the
parameterized system with local states $Q = \{\Idle, \Req, \Wait, \Crit\}$
and the local rules $t_2 = \Req \to \Wait$, $t_5 = \Crit \to \Req$, and $t_6 =
\Req \to \Idle$, together with the global rules
\[
\begin{array}{l}
t_1 = \textbf{if}\ \forall j \neq i : \{\Idle,\Req\}\ \textbf{then}\ \Idle \to \Req, \\[2pt]
t_3 = \textbf{if}\ \exists j < i : \{\Req,\Wait,\Crit\}\ \textbf{then}\ \Wait \to \Wait, \\[2pt]
t_4 = \textbf{if}\ \forall j < i : \{\Idle\}\ \textbf{then}\ \Wait \to \Crit .
\end{array}
\]
A configuration is a word over~$Q$. For instance, the configuration
$\Idle\Wait\Crit\Wait\Req$ describes an instance with five processes that are
idle, waiting, critical, waiting, and requesting, in that order. The
transition relation~$\to$ depends on the type of the transition rule. The
local rule~$t_2$ induces, for instance, the transition
$\Wait\Idle\Req\Crit\Req \to \Wait\Idle\Wait\Crit\Req$, in which the active
process changes its state from requesting to waiting. The existential
rule~$t_3$ induces the transition $\Req\Idle\Wait\Crit\Req \to
\Req\Idle\Wait\Crit\Req$, since the waiting process has a requesting process in
its left context. The same transition is not enabled from
$\Idle\Idle\Wait\Crit\Req$, since no process in the left context of the
waiting process is requesting, waiting, or critical. The universal rule~$t_4$
induces the transition $\Idle\Idle\Wait\Wait\Req \to \Idle\Idle\Crit\Wait\Req$,
since all processes in the left context of the active process are idle. From
$\Crit\Idle\Wait\Wait\Req$, neither waiting process can perform~$t_4$, since
for each of them, some process in its left context is not idle.

The initial configurations are those in which all processes are idle, such as
$\Idle\Idle$ and $\Idle\Idle\Idle\Idle\Idle$. The bad configurations are those
in which at least two processes are critical, such as $\Crit\Req\Crit$ and
$\Idle\Crit\Req\Crit\Wait\Crit$. Proving the safety property amounts to showing that no bad
configuration is reachable from an initial configuration.

\subsection{Orderings and Backward Reachability}
\label{sec:wqo}

To analyze the protocol, we use an ordering on configurations that compares
instances of different sizes, together with the theory of well-structured
transition systems~\cite{ACJT96,ACJT00,FS01}, which is presented in
Chapter~\chapref{Well-Structured Systems}. We recall here what we need.

A \emph{quasi-ordering} on a set~$C$ is a reflexive and transitive relation
${\sqsubseteq} \subseteq C \times C$. It is a \emph{well-quasi-ordering} (WQO)
if every infinite sequence $c_0, c_1, c_2, \ldots$ of elements of~$C$ contains
indices $i < j$ with $c_i \sqsubseteq c_j$. For words $u, v \in Q^*$, we write $u
\preceq v$ if $u$ is a \emph{subword} of~$v$, that is, if $u = v[i_1] \cdots
v[i_k]$ for some positions $i_1 < \cdots < i_k$. For example, $\Wait\Crit
\preceq \Req\Wait\Idle\Crit\Wait$. By Higman's lemma~\cite{Hig52}, the subword
ordering is a WQO.

For a set $X \subseteq C$, we define its \emph{upward closure} $\upc{X} = \{ c
\mid c' \sqsubseteq c \text{ for some } c' \in X \}$, and we say that $X$ is
\emph{upward-closed} if $X = \upc{X}$. The \emph{generator}
$\mathit{gen}(U)$ of an upward-closed set~$U$ is its set of minimal elements.
If $\sqsubseteq$ is a WQO, then $\mathit{gen}(U)$ is finite \lukas{as clear from the definition ...} and $U =
\upc{\mathit{gen}(U)}$, so upward-closed sets can be represented by finite
sets. Moreover, every increasing sequence $U_0 \subseteq U_1 \subseteq \cdots$
of upward-closed sets stabilizes after finitely many steps. In our example, the
set of bad configurations is the upward-closed set $\upc{\{\Crit\Crit\}}$ of
configurations that contain $\Crit\Crit$ as a subword.

A transition relation~$\to$ on~$C$ is \emph{monotonic} with respect
to~$\sqsubseteq$ if, whenever $c_1 \to c_2$ and $c_1 \sqsubseteq c_3$, there is
a configuration $c_4$ with $c_3 \to c_4$ and $c_2 \sqsubseteq c_4$. A
\emph{well-structured transition system} is a transition system with a WQO on
its configurations with respect to which its transition relation is
monotonic. If $\to$ is monotonic and $U$ is upward-closed, then $\Pre(U)$ is
upward-closed.\lukas{not hard to show}

Suppose that $\Bad$ is upward-closed. Procedure~\ref{alg:backward} computes a
finite set~$U$ with $\upc{U} = \Pre^*(\Bad)$, the set of configurations from
which $\Bad$ is reachable. It starts from the generator of~$\Bad$ and adds the
minimal predecessors of~$\upc{U}$ until no new configuration is found. A
configuration~$c$ is discarded if $c' \sqsubseteq c$ for some $c'$ already
in~$U$; we say that $c$ is \emph{subsumed} by~$c'$, since $\upc{c} \subseteq
\upc{c'}$. The system is safe if and only if $\upc{U} \cap \Init = \emptyset$.

\begin{algorithm}[ht]
\caption{Backward reachability analysis}
\label{alg:backward}
\begin{algorithmic}[1]
\REQUIRE A well-structured transition system, a set $\Init$, and the generator $G$ of an upward-closed set $\Bad$
\ENSURE Whether $\Bad$ is reachable from $\Init$
\STATE $U \leftarrow G$
\REPEAT
  \STATE $N \leftarrow \{ c \in \mathit{gen}(\Pre(\upc{U})) \mid c' \not\sqsubseteq c \text{ for all } c' \in U \}$
  \STATE $U \leftarrow U \cup N$
\UNTIL{$N = \emptyset$}
\RETURN whether $\upc{U} \cap \Init \neq \emptyset$
\end{algorithmic}
\end{algorithm}

\begin{theorem}[\cite{ACJT96,ACJT00}]
\label{thm:backward}
Procedure~\ref{alg:backward} terminates and returns the correct answer for every
transition system $(C, \to)$ with an ordering~$\sqsubseteq$ that satisfies the
following conditions:
\begin{enumerate}[nosep]
\item $\to$ is monotonic with respect to~$\sqsubseteq$;
\item $\sqsubseteq$ is a WQO;
\item for every $c \in C$, the finite set $\mathit{gen}(\Pre(\upc{\{c\}}))$ is
computable;
\item $\sqsubseteq$ is decidable;
\item for every finite set $B \subseteq C$, the emptiness of $\upc{B} \cap
\Init$ is decidable.
\end{enumerate}
\end{theorem}

Conditions~1 and~2 ensure that each set computed by the procedure is
upward-closed and that the sequence of sets stabilizes; conditions 3 to~5 make
each step effective. Petri nets, broadcast protocols, and lossy channel systems
satisfy these conditions~\cite{ACJT96,EFM99}. Our protocol satisfies
conditions~2 to~5 with respect to the subword ordering, but not condition~1, as
we show next.

\subsection{Over-Approximation}

We now define an abstract system that over-approximates the concrete one and
that satisfies the conditions of Theorem~\ref{thm:backward} with respect to the
subword ordering~$\preceq$. The main task is to obtain
a transition relation that is monotonic with respect to~$\preceq$. We first
explain why local and existential transitions are already monotonic, and why
universal transitions are not.

Consider the transition $\Idle\Req\Crit \to \Idle\Wait\Crit$ induced by the
local rule~$t_2$, and the larger configuration $\Idle\Wait\Idle\Req\Crit\Req$.
The larger configuration can perform the same local transition and reach
$\Idle\Wait\Idle\Wait\Crit\Req$, which is larger than $\Idle\Wait\Crit$. In
general, the active process of the smaller configuration also exists in the
larger configuration, and since a local transition neither checks nor changes
the states of the other processes, the larger configuration can perform the
transition while preserving the ordering.

Consider next the transition $\Req\Idle\Wait\Crit\Req \to
\Req\Idle\Wait\Crit\Req$ induced by the existential rule~$t_3$. The
configuration consists of the active process in state~$\Wait$, the left
context $\Req\Idle$, and the right context $\Crit\Req$, and the left context
contains a \emph{witness}, the requesting process, that enables the transition.
In the larger configuration $\Idle\Req\Idle\Wait\Crit\Req\Crit$, the left and
right contexts of the waiting process are larger than in the smaller
configuration, and in particular the left context still contains the witness.
Hence the larger configuration can perform the same transition.

Universal transitions, in contrast, are not monotonic. Consider the transition
$\Idle\Idle\Wait\Wait\Req \to \Idle\Idle\Crit\Wait\Req$ induced by the
universal rule~$t_4$, which is enabled since the left context of the active
process is idle. The larger configuration $\Idle\Req\Idle\Req\Wait\Wait\Req$
cannot perform the transition, since the left context of the active process
now contains requesting processes, which violate the condition.

Since local and existential transitions are monotonic, they need not be
approximated. For universal transitions, we define an abstract transition
relation~$\leadsto$ as follows. An abstract transition first deletes all
processes that violate the condition of the universal rule, and then performs
the concrete transition. For instance, we have
\[
\Idle\Req\Idle\Req\Wait\Wait\Req \leadsto \Idle\Idle\Crit\Wait\Req ,
\]
since we can first delete the two requesting processes in the left context,
obtaining $\Idle\Idle\Wait\Wait\Req$, and then perform the concrete
transition~$t_4$. For local and existential rules, $\leadsto$ coincides
with~$\to$. Every concrete transition is also an abstract transition, so the
abstract system over-approximates the concrete one, and the abstract
transition relation is monotonic with respect to~$\preceq$.

\subsection{Backward Reachability Analysis}

The abstract system satisfies the conditions of Theorem~\ref{thm:backward}. The
subword ordering is a WQO, and it is decidable, since checking whether one
word is a subword of another is easy. The abstract transition relation is
monotonic, as explained above. For a finite set~$B$ of configurations and a
regular set~$\Init$ of initial configurations, the emptiness of $\upc{B} \cap
\Init$ amounts to the emptiness of the intersection of two regular languages.
It remains to show how to compute, for a configuration~$c$ and a rule~$t$, the
generator $\Pre(t)(c)$ of the set of configurations from which $\upc{c}$ can be
reached by one abstract transition induced by~$t$. A formal definition can be
found in~\cite{ADHR07}; here we illustrate the computation for the different
types of rules of Figure~\ref{fig:mutex}.

For a local rule, we run the transition backwards on the active process. For
instance, for the local rule~$t_5$, we have $\Pre(t_5)(\Idle\Req\Wait) =
\{\Idle\Crit\Wait\}$. Strictly speaking, the predecessors also include
configurations such as $\Idle\Req\Crit\Wait$, but these are subsumed by
$\Idle\Req\Wait$ itself, and we omit them for simplicity.

For an existential rule, there are two cases, depending on whether the
configuration contains a witness. For the rule~$t_3$, we have
$\Pre(t_3)(\Req\Wait\Crit) = \{\Req\Wait\Crit\}$, since the left context of the
waiting process contains a requesting process. In contrast,
\[
\Pre(t_3)(\Idle\Wait\Crit) = \{
\Req\Idle\Wait\Crit,\, \Idle\Req\Wait\Crit,\,
\Wait\Idle\Wait\Crit,\, \Idle\Wait\Wait\Crit,\,
\Crit\Idle\Wait\Crit,\, \Idle\Crit\Wait\Crit \},
\]
since there is no witness in the left context of the active process, and we
add a witness explicitly, in each possible state and at each possible
position in the left context. The new configurations have four processes,
while the original configuration has three. The configurations generated by
the backward analysis may thus grow, and there is no a priori bound on their
size.

For a universal rule, we check whether some process in the configuration
violates the condition. For the rule~$t_4$, we have
$\Pre(t_4)(\Idle\Req\Idle\Crit\Wait) = \emptyset$, since the left context of
the critical process, which would be the active process, contains a
requesting process. Such a process would have been deleted by the abstract
transition. In contrast, $\Pre(t_4)(\Idle\Idle\Crit\Wait) =
\{\Idle\Idle\Wait\Wait\}$, since all processes in the left context of the
active process are idle.

\begin{figure}[t]
\centering
\begin{tikzpicture}[x=2.1cm,y=0.7cm,font=\small,
    gen/.style={draw,rounded corners,inner sep=3pt},
    sub/.style={text=gray,inner sep=2pt}]
  \node at (0,1) {iteration 0};
  \node at (1,1) {iteration 1};
  \node at (2,1) {iteration 2};
  \node at (3,1) {iteration 3};
  \node[gen] (cc) at (0,0) {$\Crit\Crit$};
  \node[gen] (wc) at (1,0) {$\Wait\Crit$};
  \node[gen] (rc) at (2,0) {$\Req\Crit$};
  \node[sub] (rwc) at (2,-1.2) {$\Req\Wait\Crit$};
  \node[sub] (wwc) at (2,-1.9) {$\Wait\Wait\Crit$};
  \node[sub] (cwc) at (2,-2.6) {$\Crit\Wait\Crit$};
  \node[sub] (cc2) at (3,0) {$\Crit\Crit$};
  \foreach \n in {rwc,wwc,cwc,cc2} {
    \draw[gray] (\n.south west) -- (\n.north east);
  }
  \node[draw,dashed,gray,rounded corners,fit=(rwc) (cwc)] (box) {};
  \draw[->] (cc) -- node[above,font=\scriptsize] {$t_4$} (wc);
  \draw[->] (wc) -- node[above,font=\scriptsize] {$t_2$} (rc);
  \draw[->] (wc.south) |- node[pos=0.75,above,font=\scriptsize] {$t_3$} (box.west);
  \draw[->] (rc) -- node[above,font=\scriptsize] {$t_5$} (cc2);
\end{tikzpicture}
\caption{Backward reachability analysis for the protocol of
Figure~\ref{fig:mutex}. The framed configurations form the generator of the
set of configurations from which a bad configuration is reachable in the
abstract system. The struck-out configurations are subsumed by configurations
found earlier.}
\label{fig:mutex-backward}
\end{figure}

We can now run the backward reachability analysis of
Section~\ref{sec:wqo} on the abstract system
(Figure~\ref{fig:mutex-backward}). The analysis starts from the generator
$\Crit\Crit$ of the bad configurations. The only rule that produces a critical
process is~$t_4$. Of the two critical processes in $\Crit\Crit$, only the left
one can perform~$t_4$ backwards, since the left context of the right one
contains a process that is not idle, which would have been deleted by the
abstract transition. This gives $\Pre(t_4)(\Crit\Crit) = \{\Wait\Crit\}$. From
$\Wait\Crit$, two rules are enabled backwards, both for the waiting process:
the local rule~$t_2$ gives $\Pre(t_2)(\Wait\Crit) = \{\Req\Crit\}$, and the
existential rule~$t_3$ gives $\Pre(t_3)(\Wait\Crit) =
\{\Req\Wait\Crit, \Wait\Wait\Crit, \Crit\Wait\Crit\}$. All three
configurations produced by~$t_3$ are subsumed by~$\Wait\Crit$. From
$\Req\Crit$, the only rule enabled backwards is the local rule~$t_5$, which
gives $\Crit\Crit$, which is subsumed. The rule~$t_1$ cannot be applied
backwards, since the critical process violates its condition and would have
been deleted. At this point, no new configurations are produced, and the
analysis terminates.

The set of configurations from which a bad configuration is reachable in the
abstract system is therefore $\upc{\{\Crit\Crit, \Wait\Crit, \Req\Crit\}}$.\lukas{technically, subset of this set, as it is an overapproximation (?)}
It contains no initial configuration, since the initial configurations contain
only idle processes. Hence no bad configuration is reachable in the abstract
system, and, since $\to \subseteq \leadsto$, none is reachable in the concrete
system either. In this example, the abstraction does not lose any precision
with respect to the property: the configurations that are reachable in the
concrete system are exactly those that contain none of $\Crit\Crit$,
$\Wait\Crit$, and $\Req\Crit$ as a subword. Intuitively, a process enters the
critical section only when all processes to its left are idle, and these
processes cannot leave the idle state while it is critical. We return to this
observation in Section~\ref{sec:views}.

\subsection{Extensions}

Monotonic abstraction has been extended in several directions. In many
protocols, such as Burns' and Szymanski's algorithms, a global condition is not
checked atomically: the process inspects the other processes one at a time,
and the other processes may move during the inspection. Such non-atomic global
conditions can be handled by a refinement of the abstraction~\cite{ABDR08}.
Processes may also have numerical variables, for instance clocks or tickets,
whose values are compared using gap-order constraints~\cite{ADR09}. The method
has been applied to tree topologies, to broadcast communication and binary
rendez-vous, and to German's directory-based cache coherence
protocol~\cite{ADHR07}. Since the abstraction may introduce false positives, it
has been combined with a counterexample-guided refinement loop, which
restricts the deletions allowed in the abstract transitions when a spurious
counterexample is found~\cite{ACDHHR10}. Monotonic abstraction is also related
to the verification of programs running under weak memory models, where the
buffers of the processes raise similar questions of monotonicity, see
Section~\ref{sec:other} and Chapter~\chapref{Weak Memory Models}.
